\documentclass[11pt,a4paper]{article}
\usepackage[T1]{fontenc}
\usepackage[utf8]{inputenc}
\usepackage[polish]{babel}
\usepackage{amsmath,amssymb,amsthm,geometry,hyperref}
\geometry{margin=25mm}
\newtheorem{theorem}{Twierdzenie}
\newtheorem{corollary}{Wniosek}
\newcommand{\Tr}{\operatorname{Tr}}
\newcommand{\Rea}{\operatorname{Re}}
\newcommand{\Sym}{\operatorname{Sym}}
\newcommand{\Dtr}{D_{\mathrm{tr}}}
\title{FIN: operacyjna przeszkoda dla deterministycznego sprzężenia uczącego
oraz dokładny mechanizm konwersji koherencji w geometrię}
\author{Raport badawczy FIN --- ST8651--ST8652}
\date{9 września 2026}
\begin{document}
\maketitle

\begin{abstract}
Najważniejszym wynikiem jest dokładny kontrprzykład dla proponowanego
sprzężenia $\dot\psi=iK\psi$,
$\dot K=\eta(\Pi\Rea\psi\psi^* -\gamma K)$, gdzie $\Pi$ usuwa przekątną.
Dwa zespoły stanów czystych o tej samej początkowej macierzy gęstości,
z tym samym początkowym jądrem strict, dają różne późniejsze wyniki reguły
Borna. Różnica pojawia się w drugim rzędzie czasu dla odczytu fazowego
i w trzecim rzędzie dla zwykłej populacji węzła. Podano jawne dodatnie
oszacowania przy skończonym czasie, zachowaniu dodatnich wag i dokładnej
normalizacji stanów. Dlatego ta deterministyczna reguła nie ma wspólnego
afinicznego, a więc również CPTP, rozszerzenia zgodnego ze wszystkimi
jej przewidywaniami dla stanów czystych. Jest to istotna przeszkoda dla
jej dokładnej interpretacji kwantowej, nie dowód niemożliwości FIN jako
całości ani zakaz przybliżeń klasycznych lub wielociałowych.
Drugim wynikiem jest dokładny sektor pięciu zmiennych: dostarczona
koherencja chiralna przechodzi w trwałe odkształcenie geometrii, lecz
granica zależy od zanikającego zaburzenia i nie daje ciągłego selektora.
Raport jest źródłem TeX; nie wygenerowano PDF.
\end{abstract}

\section{Zakres, pochodzenie równań i znaczenie wyniku}

Punktem wyjścia jest rzeczywista propozycja z rozdziału 4.1 dokumentu
\texttt{The FIN Kernel as an Unknown Mat.md}: układ stanu czystego
z rzutowanym źródłem Hebbowskim. Nie zastępujemy go inną teorią, aby
uzyskać kontrprzykład. Wcześniejszy audyt
\texttt{fin\_projected\_learning/REPORT.md} oddzielił ten układ od
klasycznej implementacji z zewnętrznym losowym nauczycielem, od
uśredniania cyklicznego i od jawnego rozszerzenia na stany mieszane.
Wykorzystano również wcześniejszą analizę prawa składania
\texttt{fin\_transition\_hierarchy/REPORT.md} oraz dokładne przedziały
wag w \texttt{fin\_replication\_consistency/certify.py}.
Nie twierdzimy, że ponownie udowodniono każde historyczne twierdzenie repozytorium.

Badamy
\begin{equation}\label{eq:pure}
 \dot\psi=iK\psi,\qquad
 \dot K=\eta\bigl(T(P)-\gamma K\bigr),\quad
 P=\psi\psi^*,\quad T(P)=\Pi\Rea P,
 \quad K\in\Sym_0(\mathbb R^n),\quad \eta,\gamma>0.
\end{equation}
Wszystkie czasy i parametry są bezwymiarowe; $\eta$ oznacza szybkość
uczenia, nie wykładnik tłumienia jądra. Hermitowskość $K(t)$ daje
$\|\psi(t)\|=1$ dokładnie. Nie daje jednak jednego, niezależnego od
wejścia operatora unitarnego: cała historia $K(t)$ zależy od $\psi(0)$.

W świadku strict używamy dostarczonego jądra
\begin{equation}\label{eq:strict}
 W_{ij}=\frac{\cos(0.18575d_{ij}+0.16250)}{1+d_{ij}^{9/5}},\quad
 W_{ii}=0,\qquad d_{ij}=\min(|i-j|,12-|i-j|).
\end{equation}
Nie jest to nowa derywacja parametrów strict. Główna tożsamość poniżej
obowiązuje dla dowolnego wspólnego $K(0)$, więc nie można usunąć
przeszkody przez ponowne dopasowanie tych parametrów. Dotyczy również
analogicznego sprzężenia rozpoczętego od jądra legacy, bez utożsamiania
obu jąder lub przenoszenia ich ról fizycznych.

Znaczenie wyniku jest konkretne: przed dalszą próbą uznania uczącego się
jądra za fundamentalną dynamikę kwantową trzeba rozstrzygnąć jej prawo
składania, przygotowania i pomiaru. Sama samospójność, zachowanie normy
lub lepsze dopasowanie widma tego nie załatwia. Ogólny mechanizm
zależności nieliniowej ewolucji od rozkładu zespołu jest znany z prac
Gisina i Polchinskiego \cite{Gisin1990,Polchinski1991}. Nowym wynikiem
tego audytu jest jawny świadek dla badanego równania FIN, jego kontrola
w pełnym układzie dwunastowęzłowym i certyfikaty skończonego czasu;
nie zgłaszamy pierwszeństwa dla ogólnego twierdzenia o nieliniowej QM.

\section{Dokładna przeszkoda operacyjna w oryginalnym układzie czystym}

Niech $E_j=|j\rangle\langle j|$. W płaszczyźnie węzłów $0,1$ wybierzmy
\begin{equation}\label{eq:projectors}
 P_+=\frac1{10}\begin{pmatrix}9&3\\3&1\end{pmatrix},\qquad
 P_-=\frac1{10}\begin{pmatrix}1&-3\\-3&9\end{pmatrix},\qquad
 Y=\begin{pmatrix}0&-i\\i&0\end{pmatrix}.
\end{equation}
Macierze te są osadzone w całej przestrzeni, z zerami poza wskazaną
płaszczyzną. Są to rzeczywiste, ortogonalne projektory rzędu jeden,
$P_++P_-=E_0+E_1$. Porównujemy zespoły
\begin{equation}\label{eq:ensembles}
 \mathcal E_0=\{E_0,E_1,E_2,\ldots,E_{n-1}\}/n,
 \quad
 \mathcal E_1=\{P_+,P_-,E_2,\ldots,E_{n-1}\}/n.
\end{equation}
Oba mają stan początkowy $I/n$, a w każdej realizacji ten sam $K(0)$.
Nie wprowadzamy początkowego stanu chiralnego ani zespołowo zależnego
stanu kontrolera. Każdy stan czysty ewoluuje według \eqref{eq:pure};
$\bar P_0(t),\bar P_1(t)$ oznaczają średnie po realizacjach.

\begin{theorem}[Brak wspólnego rozszerzenia afinicznego]\label{thm:affine}
Dla dowolnego $n\ge2$, wspólnego rzeczywistego symetrycznego $K(0)$
o zerowej przekątnej i $\eta>0$ zachodzi
\begin{equation}\label{eq:defect}
 \Delta P(t):=\bar P_1(t)-\bar P_0(t)
   =\frac{6\eta}{25n}\,t^2Y+O(t^3).
\end{equation}
Nie istnieje zatem wspólna afiniczna mapa macierzy gęstości, która dla
wszystkich czterech wyróżnionych stanów czystych odtwarza ten przepływ
przy wszystkich dostatecznie małych dodatnich czasach.
\end{theorem}

\begin{proof}
Mamy $\dot P=i[K,P]$ oraz
\begin{equation}\label{eq:second}
 \ddot P=i\eta[T(P),P]-i\eta\gamma[K,P]-[K,[K,P]].
\end{equation}
Dwa ostatnie składniki są liniowe w $P$ przy ustalonym wspólnym $K(0)$,
więc znikają w różnicy średnich. Dla $E_j$ mamy $T(E_j)=0$.
Bez żadnych przybliżeń,
\begin{equation}
 i[T(P_+),P_+]+i[T(P_-),P_-]=\frac{12}{25}Y.
\end{equation}
Początkowe wartości i pierwsze pochodne średnich są równe; rozwinięcie
Taylora daje \eqref{eq:defect}. Gdyby jedna mapa afiniczna $\Lambda_t$
zgadzała się na tych wejściach z badanym przepływem, równość
$P_++P_-=E_0+E_1$ wymuszałaby równość odpowiednich sum wyjściowych.
To przeczy niezerowej różnicy. Gładkość potrzebna do rozwinięcia wynika
z wielomianowego pola wektorowego; oszacowanie reszty podajemy poniżej.
\end{proof}

Odczyt $M_y=|y_+\rangle\langle y_+|$,
$|y_+\rangle=(|0\rangle+i|1\rangle)/\sqrt2$, daje
$\Tr(M_y\Delta P)=6\eta t^2/(25n)+O(t^3)$.
Nie trzeba jednak zakładać dostępności takiego odczytu fazowego.

\begin{theorem}[Różnica populacji kanonicznego węzła]\label{thm:vertex}
Przy tych samych założeniach,
\begin{equation}\label{eq:vertex}
 \Tr(E_0\Delta P(t))=-\frac{12\eta K_{01}(0)}{25n}\,t^3+O(t^4).
\end{equation}
Dla dodatniego strict jest to niezerowa różnica zwykłej populacji węzła.
\end{theorem}

\begin{proof}
Dla rzeczywistego początkowego $P$ mamy $\Rea\dot P=0$.
W przekątnej trzeciej pochodnej nieliniowa część wynosi
\[
 -\eta\{2[T(P),[K,P]]+[K,[T(P),P]]\}.
\]
Pozostałe części znoszą się między zespołami albo mają zerową przekątną.
Suma dla $P_+,P_-$ ma na przekątnej wartości
$(-72K_{01}/25,72K_{01}/25,0,\ldots,0)$, przed czynnikiem $\eta/n$.
Współczynniki sprzężeń do wszystkich pozostałych węzłów znikają:
$T(P_\pm)$ i $[T(P_\pm),P_\pm]$ są wspierane na płaszczyźnie $0,1$.
Podzielenie przez $3!$ daje \eqref{eq:vertex}. Tożsamości sprawdza
również dokładny rachunek symboliczny z dowolnymi dodatkowymi sprzężeniami.
\end{proof}

\subsection{Skończony czas, dodatnie krawędzie i kontrola reszty}

Dla $n=12$, $K(0)=W$, $\eta=1/5$, $\gamma=1/20$, $T=10^{-4}$
nie opieramy się tylko na symbolu $O(t^3)$. Dokładne przedziały wag
potwierdzają $\|W\|_F<3-\eta T$ i $W_{01}>2/5$.
Z całkowej postaci równania dla $K$ i $\|T(P)\|_F\le1$ dostajemy
$\|K\|_F\le B=3$ na całym $[0,T]$. Zdefiniujmy
\begin{align}
 L&=\eta(1+\gamma B)=23/100,&
 J&=\eta(2B+\gamma L)=12023/10000,\\
 A_2&=2L+4B^2,&
 A_3&=2J+12BL+8B^3=1133423/5000,\\
 H&=\eta(A_2+\gamma J),&
 A_4&=2H+12BJ+6LA_2+2BA_3=734156623/500000.
\end{align}
Nierówność $\|[A,B]\|_F\le2\|A\|_{\rm op}\|B\|_F$ daje kolejno
$\|K'\|_F\le L$, $\|K''\|_F\le J$,
$\|P''\|_F\le A_2$, $\|P'''\|_F\le A_3$ i
$\|P''''\|_F\le A_4$. Wszystkie cztery różniące się realizacje
spełniają te same oszacowania; pozostałe dziesięć znosi się dokładnie.

Stąd reszta dla $M_y$ nie przekracza $2A_3T^3/(3n)$, a dla $E_0$
nie przekracza $A_4T^4/(6n)$. Otrzymujemy
\begin{align}\label{eq:finite}
 \Tr(M_y\Delta P(T))
 &\ge\frac{2466577}{90000000000000000}>2.74\cdot10^{-11},\\
 \Tr(E_0\bar P_0(T))-\Tr(E_0\bar P_1(T))
 &\ge\frac{417843377}{360000000000000000000000}>10^{-15}.
\end{align}
Jednocześnie $K_{ij}(t)\ge W_{ij}-LT>0$ dla $i\ne j$ na tym samym
przedziale, co jest sprawdzane na dokładnych dolnych granicach wag.
To bardzo zachowawcze oszacowania matematyczne, nie czułość istniejącego
aparatu. Nie przeliczamy bezwymiarowego $T$ na czas laboratoryjny.

Niezależna całka numeryczna pełnego układu czystego dla czterech wejść
potwierdza znaki i rzędy. Przy $t=0.1$ daje około
$3.9683\cdot10^{-5}$ różnicy prawdopodobieństwa $M_y$ oraz
$3.7392\cdot10^{-6}$ różnicy populacji węzła. Te większe wartości są
wynikami numerycznymi, nie zamiennikami certyfikatów \eqref{eq:finite}.

\subsection{Pamięć kontrolera i granica możliwej dokładności}

Dowolne standardowe urządzenie kwantowe, z dowolną pamięcią i środowiskiem
o wspólnym stanie początkowym niezależnym od nieznanego wejścia,
wyznacza przy ustalonym czasie mapę CPTP na wejściowej macierzy gęstości.
Nie trzeba zakładać dynamiki Markowa. Twierdzenie~\ref{thm:affine}
wyklucza więc dokładną realizację \eqref{eq:pure} w tej klasie.

Jeśli błąd w odległości śladowej na każdym z czterech wejść jest najwyżej
$\epsilon$, nierówność trójkąta daje
$|\Tr(M_y\Delta P)|\le4\epsilon/n$.
Dla świadka strict i czasu $T$ przynajmniej jedno wejście musi zatem
mieć błąd co najmniej $2466577/30000000000000000$.
Nie jest to wykluczenie użytecznego przybliżenia w skończonej tolerancji.

\section{Niezależny, zamknięty świadek dwupoziomowy}

Nie jest on zastępstwem dowodu dla strict, lecz kontrolą analityczną.
Dla $K=kX$, $X=\left(\begin{smallmatrix}0&1\\1&0\end{smallmatrix}\right)$,
początkowych wektorów Blocha $\pm(a,0,b)$, $a=3/5$, $b=4/5$,
oraz $k(0)=k_0$, mamy $\lambda=\eta\gamma$ i
\begin{align}
 k_\pm(t)&=k_0e^{-\lambda t}\pm\frac{a}{2\gamma}(1-e^{-\lambda t}),\\
 \Theta_0(t)&=\frac{2k_0}{\lambda}(1-e^{-\lambda t}),&
 \delta(t)&=\frac a\gamma\left[t-\frac{1-e^{-\lambda t}}\lambda\right].
\end{align}
Zespół osi $Z$ pozostaje średnio $I/2$, podczas gdy zespół przechylony ma
średni wektor Blocha
\[
 \bar r(t)=b\sin\delta(t)(0,\cos\Theta_0(t),-\sin\Theta_0(t)).
\]
Dlatego $\Dtr=\tfrac b2|\sin\delta|$. Dla $\eta=1/5$, $\gamma=1/20$,
$k_0=t=1$ naprzemienne oszacowania szeregu wykładniczego dają
\[
 \frac{299}{5000}<\delta<\frac{119601}{2000000},\qquad
 \Dtr>\frac{239}{10000},\qquad k_-(1)>\frac{93}{100}.
\]
Użyto $\sin\delta\ge\delta-\delta^3/6$. Każda wspólna mapa CPTP musi
na co najmniej jednym z czterech wejść popełnić błąd odległości śladowej
większy niż $239/20000$. W tym świadku przeszkoda nie jest wyłącznie
formalnym efektem nieskończenie małego czasu. Numeryczna kontrola
oryginalnego równania zgadza się z zamkniętym wzorem do około $2.3\cdot10^{-15}$.

\section{Próby falsyfikacji i rzeczywiste drogi obejścia}

\paragraph{Tłumienie, kształt jądra i dodatniość.}
Współczynnik \eqref{eq:defect} nie zawiera $K(0)$ ani $\gamma$.
Zmiana tych danych nie usuwa przeszkody. Dodatniość została opłacona
osobno; świadek nie korzysta z wcześniejszego problemu ujemnych krawędzi.
Przy $\eta=0$ znika i odzyskujemy wspólną ewolucję unitarną, co
potwierdza test kontrolny. Szczególne zespoły o $a=0$ lub $b=0$ też
nie dają tego kontrastu; nie każda para rozkładów musi go wykazywać.

\paragraph{Propagacja przez laplasjan.}
Zastąpienie $iK\psi$ przez $-iA(K)\psi$, gdzie
$A(K)=\operatorname{diag}(K\mathbf1)-K$, nie usuwa drugiego rzędu.
W nieliniowej części pojawia się $T(P)-\operatorname{diag}(T(P)\mathbf1)$.
Dla projektorów \eqref{eq:projectors} dodana przekątna jest skalarem
na wspierającej je płaszczyźnie, więc z nimi komutuje. Współczynnik
\eqref{eq:defect} pozostaje dokładnie taki sam. Nie przenosimy na ten
zmieniony układ bez ponownego rachunku wszystkich stałych reszty z sekcji 2.

\paragraph{Równanie zapisane bezpośrednio dla macierzy mieszanej.}
Układ $\dot\rho=i[K,\rho]$, $\dot K=\eta(T(\rho)-\gamma K)$ jest
dobrze określonym innym przepisem. Dla $\rho(0)=I/n$ daje
$\rho(t)=I/n$, $K(t)=e^{-\eta\gamma t}K(0)$. Nie zgadza się jednak
automatycznie ze średnią po realizacjach stanów czystych.
Nazwanie $\rho$ macierzą gęstości nie dowodzi zgodności tych dwóch operacji.

\paragraph{Dokładny zakres wniosku o sygnalizacji.}
Jeżeli przyjmujemy standardowe przygotowanie zdalne z maksymalnie
splątanego stanu, regułę Borna i ewolucję warunkowo przygotowanych
stanów według \eqref{eq:pure}, wybór zdalnej bazy przygotowuje
\eqref{eq:ensembles} bez zmiany lokalnego $I/n$. Różne późniejsze
populacje dawałyby w tym złożeniu przewidywanie sygnalizacji.
Jest to wniosek warunkowy, nie zaobserwowany sygnał i nie dowód, że
sama abstrakcyjna zasada braku sygnalizacji zakazuje każdej nieliniowości.
Założenia o zespołach są istotne również w literaturze
\cite{Simon2001,Bona2003}; tutaj zostały przyjęte jawnie, a nie wyprowadzone
z samego zachowania normy.

\paragraph{Model klasyczny albo wielociałowy.}
Przepływ na rzeczywistych stanach ontologicznych $(\psi,K)$ i na
rozkładach prawdopodobieństwa tych stanów jest matematycznie dopuszczalny.
Wtedy sama średnia $\rho$ nie musi być pełnym stanem fizycznym.
Przybliżenie pola średniego także nie jest wykluczone: dla wielu
jednakowo przygotowanych kopii wejścia
$\sum_jp_jP_j^{\otimes N}$ mogą się różnić mimo równości
$\sum_jp_jP_j$. Dla dwóch kopii naszych par
\[
 \Tr\!\left[(X\otimes Z)
 \left\{\frac{P_+\otimes P_++P_-\otimes P_-}{2}
       -\frac{E_0\otimes E_0+E_1\otimes E_1}{2}\right\}\right]
 =\frac{12}{25}.
\]
Nieliniowość może więc odzwierciedlać dodatkowe korelacje wejściowe,
a nie lokalną operację na jednym nieznanym stanie kwantowym. To wiąże
wynik z wcześniejszą przeszkodą domknięcia hierarchii: potrzebny jest
rzeczywisty dostawca wyższych momentów, kontrolera i prawa składania.
Zaszumione instrumenty kwantowe lub autonomiczna teoria wielociałowa
mogą być inną drogą, ale ich dodanie nie jest dowodem, że FIN już je wyprowadza.

\section{Drugi wynik: dokładna konwersja koherencji w geometrię}

Ten wynik dotyczy jawnego rozszerzenia mieszanego, nie usuwa przeszkody
operacyjnej powyżej. Dla $i,j=0,\ldots,11$ połóżmy
\[
 E_{ij}=\tfrac13\cos(\pi(i+j)/6),\qquad
 F_{ij}=\tfrac13\sin(\pi(i+j)/6)
 \quad\text{gdy $i+j$ jest nieparzyste},
\]
a zero w pozostałych przypadkach, oraz $G=[E,F]/(2i)$.
Dokładna algebra nad $\mathbb Q(\sqrt3,i)$ daje
\[
 [E,F]=2iG,\quad[F,G]=2iE,\quad[G,E]=2iF,\quad
 E^2=F^2=G^2=P_1+P_5.
\]
Antykomutatory parami są zerowe. Wszystkie trzy macierze mają zerową
przekątną i sumy wierszy oraz komutują z dowolnym rzeczywistym
symetrycznym jądrem cyrkulantowym $C_{12}$.

Dla $R>0$, $r_*=R/\gamma$, $\kappa=\eta\gamma^2/R$, $\tau=r_*t$,
podstawienie
\[
 K=W+r_*(xE+yF),\qquad
 \rho=I/12+\gamma W+R(uE+vF+zG)
\]
daje dokładny, niezmienniczy sektor pełnego równania:
\begin{equation}\label{eq:spin}
 x'=\kappa(u-x),\quad y'=\kappa(v-y),\quad
 u'=-2yz,\quad v'=2xz,\quad z'=-2(xv-yu).
\end{equation}
Nie narzucamy tego rzutu podczas ewolucji. Jest to ograniczenie danych
początkowych, a nie dodatkowy człon dynamiki.

Niech $w=x+iy$, $b=u+iv$. Sfera $|b|^2+z^2=1$ jest niezmiennicza,
a $|w(0)|\le1$ pociąga $|w(\tau)|\le1$.
Dopuszczalność stanu zapewnia $R\le\min(c_1,c_5)$,
$c_j=1/12+\gamma\lambda_j(W)$, przy nieujemnych pozostałych $c_j$.
Dodatniość krawędzi zapewnia
$r_*/3<\min(W_1,W_3,W_5)$. Przykład
$\gamma=1/20$, $R=1/2000$, $\eta=1/5$ spełnia te nierówności według
dokładnych przedziałów strict. Populacje pozostają dokładnie $1/12$.

\begin{theorem}[Zbieżność, lecz nie ciągły wybór geometrii]
Dla $w(0)=0$, $b(0)=\epsilon e^{i\phi}$,
$z(0)=\sqrt{1-\epsilon^2}$, $0<\epsilon\le1$, rozwiązanie
\eqref{eq:spin} zbiega do jednego punktu
$w=b=e^{i\Theta}$, $z=0$. Zbiór wszystkich punktów skupienia mapy
końcowej, gdy dane początkowe dążą do bieguna $w=b=0,z=1$, jest całym
okręgiem tych geometrii. Nie istnieje ciągły selektor w tym biegunie.
\end{theorem}

\begin{proof}
Na zwartej przestrzeni rozważanych trajektorii funkcja
\[
 V=|w|^2-2\Rea(\bar wb),\quad
 V+1=|w-b|^2+z^2,\quad V'=-2\kappa|w-b|^2
\]
ma największy zbiór niezmienniczy zerowej pochodnej złożony z dwóch
biegunów i okręgu równowag. Dla $q=|w|^2$ zachodzi
$q''+\kappa q'=2|w'|^2$ oraz $q'(0)=0$.
Ponieważ $w'(0)\ne0$, mamy $q'>0$ dla każdego dodatniego czasu.
Zasada niezmienniczości wyklucza bieguny jako granice, dając
$q\to1$, $z\to0$, $w-b\to0$.

Aby wykluczyć wędrówkę po okręgu, od dowolnego $\tau_0>0$ wybierzmy
ciągłą fazę $\theta=\arg w$. Mamy $\theta'=-\kappa z'/(2q)$, a więc
\[
 \theta(T)-\theta(\tau_0)
 =-\frac\kappa2[z/q]_{\tau_0}^{T}
  -\frac\kappa2\int_{\tau_0}^{T}\frac{zq'}{q^2}\,d\tau.
\]
Całka jest bezwzględnie zbieżna, bo $|z|\le1$, $q'>0$ i
$\int q'/q^2=1/q(\tau_0)-1/q(T)$. Istnieje więc jedna końcowa faza.
Symetria $(w,b)\mapsto e^{i\phi}(w,b)$ obraca tę fazę o $\phi$.
Dla dowolnie małego $\epsilon$ można otrzymać dowolny punkt okręgu;
dla dokładnie zerowego zaburzenia biegun pozostaje stacjonarny.
\end{proof}

Linearizacja w dodatnim biegunie ma niestabilny pierwiastek
$\mu=(-\kappa+\sqrt{\kappa^2+8i\kappa})/2$, $\Rea\mu>0$.
Końcowe odkształcenie ma normę Frobeniusa $2R/\gamma$ niezależną od
amplitudy niezerowego zaburzenia, lecz czas jego rozwoju rośnie, gdy
zaburzenie maleje. Nie ma nieciągłości rozwiązań przy ustalonym skończonym
czasie. Widmo stanu, jego entropia von Neumanna i wszystkie populacje
początkowe są te same dla tych danych; informacja rozstrzygająca leży
w koherencji. Źródło tej informacji i początkowego strict nadal jest dane.

Zaobserwowane logarytmiczne nawijanie fazy dla $K(0)=W$ ma zgodny
współczynnik $-\Im\mu/\Rea\mu$, ale pełna asymptotyka tej konkretnej
krzywej wejściowej pozostaje nieudowodniona. Nie jest przesłanką
powyższego twierdzenia. Nie dowodzimy również stabilności w całej
przestrzeni dwunastowęzłowej ani nowego uniwersalnego prawa skali.

\section{Bilans i najbardziej rozstrzygający dalszy kierunek}

\paragraph{Udowodniono.}
Oryginalne czyste sprzężenie \eqref{eq:pure} ma jawną zależność wyniku
od rozkładu zespołu, a zatem nie ma wspólnego afinicznego rozszerzenia
na standardowe stany kwantowe. Przeszkoda występuje dla strict przy
dodatnich krawędziach, obejmuje kanoniczne odczyty i nie znika przez
zmianę tłumienia ani przejście do propagatora laplasjanowego.
Osobno udowodniono dokładny sektor konwersji koherencji i brak
ciągłego, niezależnego od zaburzeń wyboru geometrii w tym sektorze.

\paragraph{Obalono w określonym zakresie.}
Nie można utożsamiać zachowania normy każdej realizacji ze wspólną
operacją unitarną, ani traktować zapisu równania dla średniej macierzy
gęstości jako automatycznie zgodnego ze średnią po przygotowaniach.
Nie można przypisać temu sprzężeniu bezwarunkowej selekcji strict
niezależnej od wejściowych danych koherencyjnych.

\paragraph{Pozostaje warunkowe lub otwarte.}
Wniosek o sygnalizacji wymaga jawnego standardowego prawa przygotowania
i ewolucji stanów warunkowych. Modele klasyczne, rozkłady na stanach
ontologicznych, przybliżenia pola średniego i inne prawa złożone nie
zostały wykluczone. Nie ma tu laboratoryjnego testu FIN, źródła jednostek,
wyprowadzenia stanu początkowego, pełnego mostu legacy--strict, transferu
ról fizycznych, rozstrzygnięcia QW-2191, SM/GR, $L_{\rm total}$ ani ToE.
Nadsoliton pozostaje pierwotną informacją w stanie solitonicznym;
nie dodajemy pod nim osobnej warstwy informacyjnej.

\paragraph{Następny test o największej wartości informacyjnej.}
Należy wskazać jedno rzeczywiście wyprowadzone z FIN, skończone prawo
mikroskopowe dla układu i jego kontrolera lub otoczenia, określić prawo
składania i wykazać zgodność operacyjną z przygotowaniem i pomiarem.
Jeżeli ma ono dawać \eqref{eq:pure} tylko jako przybliżenie wielociałowe,
potrzebna jest kontrolowana granica z błędem, fluktuacjami i jawnym
źródłem wyższych korelacji, nie założenie ukryte pod słowem
``samouczenie''. Dopiero wtedy warto badać, czy strict wynika z tego
prawa bez zakodowania celu w danych wejściowych. To bardziej
rozstrzygające od kolejnego dopasowania jądra, ponieważ testuje możliwość
przejścia od poprawnej matematyki do jednej spójnej fizyki.

\section*{Reprodukowalność i kontrola zakresu}

Pliki \texttt{fin\_chiral\_selection/research.py} i
\texttt{operational.py} odtwarzają obliczenia; dwa pliki testowe
sprawdzają dokładną algebrę oraz niezależne kontrole numeryczne.
Dokładne certyfikaty opierają się na liczbach wymiernych, symbolicznych
komutatorach i zewnętrznie zaokrąglonych przedziałach wag, nie na
małych resztach zmiennoprzecinkowych. Testy obejmują zerowe uczenie,
szczególne niewrażliwe zespoły, propagację laplasjanową, normę stanów,
dodatniość, zmianę chiralności i fazy oraz niezależny wzór dwupoziomowy.
Plik \texttt{verification.json} zapisuje rzeczywiście wykonany replay,
wersje narzędzi i identyczność źródeł; nie zastępuje powyższych dowodów.
Nie instalowano ciężkiego systemu algebraicznego, nie uruchamiano
zatrzymanej kampanii P475/P485/P487 i nie generowano PDF.

\begin{thebibliography}{9}
\bibitem{Gisin1990}
N.~Gisin, \emph{Weinberg's non-linear quantum mechanics and supraluminal
communications}, Physics Letters A 143 (1990), 1--2.
\href{https://doi.org/10.1016/0375-9601(90)90786-N}{DOI: 10.1016/0375-9601(90)90786-N}.
\bibitem{Polchinski1991}
J.~Polchinski, \emph{Weinberg's nonlinear quantum mechanics and the
Einstein-Podolsky-Rosen paradox}, Physical Review Letters 66 (1991), 397.
\href{https://doi.org/10.1103/PhysRevLett.66.397}{DOI: 10.1103/PhysRevLett.66.397}.
\bibitem{Simon2001}
C.~Simon, V.~Bužek, N.~Gisin, \emph{The no-signaling condition and quantum
dynamics}, Physical Review Letters 87 (2001), 170405.
\href{https://arxiv.org/abs/quant-ph/0102125}{arXiv: quant-ph/0102125}.
\bibitem{Bona2003}
P.~Bóna, \emph{Comment on ``No-Signaling Condition and Quantum Dynamics''},
Physical Review Letters 90 (2003), 208901.
\href{https://arxiv.org/abs/quant-ph/0201002}{arXiv: quant-ph/0201002}.
\end{thebibliography}
\end{document}
